\documentclass[11pt]{article}

% Shared notes settings for Elements of AIML lectures (UPES).
% Keep this file free of \documentclass to allow reuse via \input.

\usepackage[utf8]{inputenc}
\usepackage[T1]{fontenc}
\usepackage[a4paper,margin=1in]{geometry}
\usepackage{hyperref}
\usepackage{xurl}
\usepackage{booktabs}
\usepackage{amsmath}
\usepackage{amssymb}
\usepackage{listings}
\usepackage{xcolor}
\usepackage{enumitem}
\usepackage{parskip}

% Code listings (general-purpose).
\lstset{
  basicstyle=\ttfamily\small,
  keywordstyle=\color{blue},
  commentstyle=\color{gray},
  stringstyle=\color{teal},
  showstringspaces=false,
  columns=fullflexible,
  frame=single,
  framerule=0.2pt,
  breaklines=true
}

\setlist[itemize]{topsep=0.3em,itemsep=0.2em}
\setlist[enumerate]{topsep=0.3em,itemsep=0.2em}

% Simple per-section source line.
\newcommand{\notesources}[1]{\par\noindent{\small\color{gray}\textit{Sources:} \sloppy #1\par}}

% Unit-wise source strings for this subject (used by slides/notes).
% Path is relative to the lecture `latex/` folder (the compilation CWD).
\IfFileExists{../../../_shared/aiml-sources.tex}{%
  % Source strings for Elements of AIML (used by slides + notes).

\newcommand{\AIMLRepoURL}{https://github.com/tali7c/Elements-of-AIML}

\newcommand{\AIMLLabSources}{%
Provost and Fawcett, \textit{Data Science for Business}.;%
McKinney, \textit{Python for Data Analysis}.;%
WEKA Documentation (Waikato Environment for Knowledge Analysis), accessed 2026-02-28.;%
Matplotlib and Seaborn documentation, accessed 2026-02-28.%
}

\newcommand{\AIMLUnitISources}{%
Rich and Knight, \textit{Artificial Intelligence}, McGraw Hill, 2017.;%
Russell and Norvig, \textit{Artificial Intelligence: A Modern Approach} (4e), Ch. 1--2 (Introduction, Intelligent Agents).;%
Poole and Mackworth, \textit{Artificial Intelligence: Foundations of Computational Agents} (2e), selected sections.%
}

\newcommand{\AIMLUnitIISources}{%
Russell and Norvig, \textit{Artificial Intelligence: A Modern Approach} (4e), Ch. 7--9 (Logical Agents, First-Order Logic, Inference).;%
Huth and Ryan, \textit{Logic in Computer Science} (2e), selected sections on propositional and first-order logic.;%
Genesereth and Nilsson, \textit{Logical Foundations of Artificial Intelligence}, selected chapters.%
}

\newcommand{\AIMLUnitIIISources}{%
Mueller and Massaron, \textit{Machine Learning for Dummies}, For Dummies, 2016.;%
Mitchell, \textit{Machine Learning}, selected chapters on learning basics and evaluation.;%
Scikit-learn documentation, accessed 2026-02-28.%
}

\newcommand{\AIMLUnitIVSources}{%
Mueller and Massaron, \textit{Machine Learning for Dummies}, For Dummies, 2016.;%
James, Witten, Hastie, and Tibshirani, \textit{An Introduction to Statistical Learning} (2e), selected chapters.;%
Scikit-learn documentation, accessed 2026-02-28.%
}

\newcommand{\AIMLUnitVSources}{%
NIST, \textit{AI Risk Management Framework (AI RMF 1.0)}, 2023.;%
Barocas, Hardt, and Narayanan, \textit{Fairness and Machine Learning} (online book), selected sections.;%
Knaflic, \textit{Storytelling with Data}, selected chapters on communicating results.%
}
%
}{%
  % Faculty copies live under `private/` so their compile CWD is different.
  \IfFileExists{../../../../public/_shared/aiml-sources.tex}{%
    % Source strings for Elements of AIML (used by slides + notes).

\newcommand{\AIMLRepoURL}{https://github.com/tali7c/Elements-of-AIML}

\newcommand{\AIMLLabSources}{%
Provost and Fawcett, \textit{Data Science for Business}.;%
McKinney, \textit{Python for Data Analysis}.;%
WEKA Documentation (Waikato Environment for Knowledge Analysis), accessed 2026-02-28.;%
Matplotlib and Seaborn documentation, accessed 2026-02-28.%
}

\newcommand{\AIMLUnitISources}{%
Rich and Knight, \textit{Artificial Intelligence}, McGraw Hill, 2017.;%
Russell and Norvig, \textit{Artificial Intelligence: A Modern Approach} (4e), Ch. 1--2 (Introduction, Intelligent Agents).;%
Poole and Mackworth, \textit{Artificial Intelligence: Foundations of Computational Agents} (2e), selected sections.%
}

\newcommand{\AIMLUnitIISources}{%
Russell and Norvig, \textit{Artificial Intelligence: A Modern Approach} (4e), Ch. 7--9 (Logical Agents, First-Order Logic, Inference).;%
Huth and Ryan, \textit{Logic in Computer Science} (2e), selected sections on propositional and first-order logic.;%
Genesereth and Nilsson, \textit{Logical Foundations of Artificial Intelligence}, selected chapters.%
}

\newcommand{\AIMLUnitIIISources}{%
Mueller and Massaron, \textit{Machine Learning for Dummies}, For Dummies, 2016.;%
Mitchell, \textit{Machine Learning}, selected chapters on learning basics and evaluation.;%
Scikit-learn documentation, accessed 2026-02-28.%
}

\newcommand{\AIMLUnitIVSources}{%
Mueller and Massaron, \textit{Machine Learning for Dummies}, For Dummies, 2016.;%
James, Witten, Hastie, and Tibshirani, \textit{An Introduction to Statistical Learning} (2e), selected chapters.;%
Scikit-learn documentation, accessed 2026-02-28.%
}

\newcommand{\AIMLUnitVSources}{%
NIST, \textit{AI Risk Management Framework (AI RMF 1.0)}, 2023.;%
Barocas, Hardt, and Narayanan, \textit{Fairness and Machine Learning} (online book), selected sections.;%
Knaflic, \textit{Storytelling with Data}, selected chapters on communicating results.%
}
%
  }{%
    % Source strings for Elements of AIML (used by slides + notes).

\newcommand{\AIMLRepoURL}{https://github.com/tali7c/Elements-of-AIML}

\newcommand{\AIMLLabSources}{%
Provost and Fawcett, \textit{Data Science for Business}.;%
McKinney, \textit{Python for Data Analysis}.;%
WEKA Documentation (Waikato Environment for Knowledge Analysis), accessed 2026-02-28.;%
Matplotlib and Seaborn documentation, accessed 2026-02-28.%
}

\newcommand{\AIMLUnitISources}{%
Rich and Knight, \textit{Artificial Intelligence}, McGraw Hill, 2017.;%
Russell and Norvig, \textit{Artificial Intelligence: A Modern Approach} (4e), Ch. 1--2 (Introduction, Intelligent Agents).;%
Poole and Mackworth, \textit{Artificial Intelligence: Foundations of Computational Agents} (2e), selected sections.%
}

\newcommand{\AIMLUnitIISources}{%
Russell and Norvig, \textit{Artificial Intelligence: A Modern Approach} (4e), Ch. 7--9 (Logical Agents, First-Order Logic, Inference).;%
Huth and Ryan, \textit{Logic in Computer Science} (2e), selected sections on propositional and first-order logic.;%
Genesereth and Nilsson, \textit{Logical Foundations of Artificial Intelligence}, selected chapters.%
}

\newcommand{\AIMLUnitIIISources}{%
Mueller and Massaron, \textit{Machine Learning for Dummies}, For Dummies, 2016.;%
Mitchell, \textit{Machine Learning}, selected chapters on learning basics and evaluation.;%
Scikit-learn documentation, accessed 2026-02-28.%
}

\newcommand{\AIMLUnitIVSources}{%
Mueller and Massaron, \textit{Machine Learning for Dummies}, For Dummies, 2016.;%
James, Witten, Hastie, and Tibshirani, \textit{An Introduction to Statistical Learning} (2e), selected chapters.;%
Scikit-learn documentation, accessed 2026-02-28.%
}

\newcommand{\AIMLUnitVSources}{%
NIST, \textit{AI Risk Management Framework (AI RMF 1.0)}, 2023.;%
Barocas, Hardt, and Narayanan, \textit{Fairness and Machine Learning} (online book), selected sections.;%
Knaflic, \textit{Storytelling with Data}, selected chapters on communicating results.%
}
%
  }%
}


\title{Elements of AIML\\Unit 04 -- Lecture 02 Notes\\Supervised Learning (Classification)}
\author{Tofik Ali}
\date{\today}

\begin{document}
\maketitle
\tableofcontents

\section{Lecture Overview}
This lecture covers \textbf{classification} and the evaluation tools needed to judge classification models correctly, especially under class imbalance.
\notesources{\AIMLUnitIVSources}

\section{How to use these notes (self-study)}
Start with the confusion matrix and be sure you can explain TP/FP/TN/FN in a real scenario (medical test, fraud detection).
Then learn how each metric answers a different question (``How many alarms are correct?'' vs ``How many true cases did we catch?'').
Finally, study thresholds and imbalance, because most real classification problems have unequal error costs.

\section{Learning Outcomes}
\begin{itemize}[leftmargin=*]
  \item Define classification and give examples.
  \item Understand and use confusion matrix concepts.
  \item Compute and interpret accuracy, precision, recall, and F1.
  \item Explain threshold effects and why imbalance changes metric choice.
\end{itemize}

\section{Keywords}
\begin{itemize}[leftmargin=*]
  \item \textbf{Classification:} predicting discrete labels.
  \item \textbf{Confusion matrix:} TP/FP/TN/FN table.
  \item \textbf{Precision:} correctness among predicted positives.
  \item \textbf{Recall:} coverage of true positives.
  \item \textbf{Threshold:} rule to convert score into a class.
  \item \textbf{Imbalanced data:} class frequencies are highly unequal.
\end{itemize}

\section{Abbreviations and symbols}
\begin{itemize}[leftmargin=*]
  \item TP/FP/TN/FN: true/false positives/negatives.
  \item Precision $= TP/(TP+FP)$,\; Recall/TPR $= TP/(TP+FN)$.
  \item FPR $= FP/(FP+TN)$,\; F1: harmonic mean of precision and recall.
  \item ROC-AUC: area under ROC curve; PR-AUC: area under precision-recall curve.
  \item Threshold $\tau$: score $\ge \tau \Rightarrow$ predict positive.
\end{itemize}

\section{Classification basics}
Classification predicts a label such as spam/not-spam. Models may output:
\begin{itemize}[leftmargin=*]
  \item a hard label directly, or
  \item a score/probability, then a threshold is applied.
\end{itemize}

\subsection{Common classifiers (high-level overview)}
\begin{itemize}[leftmargin=*]
  \item \textbf{Logistic regression}: linear model that outputs probabilities for classes.
  \item \textbf{k-Nearest Neighbors (kNN)}: predicts based on majority vote of nearest samples.
  \item \textbf{Decision trees / Random forests}: learn if-then style splits; handle non-linearity well.
  \item \textbf{Naive Bayes}: fast probabilistic classifier; often strong on text features.
  \item \textbf{SVM}: margin-based classifier; powerful with the right kernel/features.
\end{itemize}
This course focuses on \textbf{evaluation concepts} rather than algorithm details.

\section{Confusion matrix}
For binary classification:
\begin{center}
\begin{tabular}{c|cc}
  & Pred + & Pred - \\\hline
  True + & TP & FN \\
  True - & FP & TN \\
\end{tabular}
\end{center}
\begin{itemize}[leftmargin=*]
  \item TP: correctly predicted positive.
  \item FP: predicted positive but actually negative.
  \item FN: predicted negative but actually positive (miss).
  \item TN: correctly predicted negative.
\end{itemize}

\subsection{Worked example (interpretation)}
Suppose a disease test model produced: TP=40, FP=10, FN=20, TN=930.
\begin{itemize}[leftmargin=*]
  \item Many TNs make accuracy look high.
  \item FN means \textbf{missed sick} people (can be costly).
  \item FP means \textbf{false alarm} (unnecessary follow-up tests).
\end{itemize}

\section{Metrics}
\subsection{Accuracy}
\[
  \text{Accuracy} = \frac{TP + TN}{TP+TN+FP+FN}
\]

\subsection{Precision, recall, F1}
\[
  \text{Precision} = \frac{TP}{TP+FP},\qquad
  \text{Recall} = \frac{TP}{TP+FN}
\]
F1 score:
\[
  F1 = 2\cdot \frac{\text{Precision}\cdot\text{Recall}}{\text{Precision}+\text{Recall}}
\]

\subsection{Mini computation (from the worked confusion matrix)}
For TP=40, FP=10, FN=20:
\[
  \text{Precision}=\frac{40}{50}=0.80,\qquad
  \text{Recall}=\frac{40}{60}\approx 0.67
\]
\[
  F1 \approx 2\cdot \frac{0.80\cdot 0.67}{0.80+0.67}\approx 0.73
\]

\subsection{Interpretation}
\begin{itemize}[leftmargin=*]
  \item High precision: few false alarms.
  \item High recall: few missed positives.
  \item Trade-off depends on domain costs (medical diagnosis vs spam filtering).
\end{itemize}

\section{Thresholds and ROC-AUC (intuition)}
If a model outputs probabilities, the threshold controls the precision/recall trade-off.
ROC curve plots:
\begin{itemize}[leftmargin=*]
  \item TPR (same as recall) vs FPR ($FP/(FP+TN)$) across thresholds.
\end{itemize}
AUC summarizes ranking quality (probability that a random positive is ranked higher than a random negative).

\subsection{ROC vs PR curves (important note)}
\begin{itemize}[leftmargin=*]
  \item ROC curves consider TPR vs FPR and can look optimistic when negatives dominate.
  \item Precision-Recall curves focus on the positive class and are often more informative for rare-event detection.
\end{itemize}

\subsection{Choosing a threshold (practical mindset)}
Choose $\tau$ based on the \textbf{cost of errors}:
\begin{itemize}[leftmargin=*]
  \item If FN is very costly (disease screening), lower $\tau$ to increase recall.
  \item If FP is very costly (blocking legitimate payments), raise $\tau$ to increase precision.
\end{itemize}

\section{Imbalanced data}
When one class is rare, accuracy can be very misleading.
Example: if 99\% are negative, predicting always negative yields 99\% accuracy but 0 recall for positives.

\subsection{What to do?}
\begin{itemize}[leftmargin=*]
  \item Use better metrics (precision/recall/F1, PR-AUC).
  \item Use resampling or class weights (covered in lab Experiment-09).
  \item Tune threshold based on costs.
\end{itemize}

\section{Practice (with answers)}
\subsection*{P1. In fraud detection, which is usually worse: FP or FN?}
\textbf{Answer:} FN is often worse because it misses fraud. However, FP can also be costly due to blocking real customers; the correct answer depends on business cost.

\subsection*{P2. If a model has very high precision but very low recall, what does it mean?}
\textbf{Answer:} it predicts positive rarely, but when it does, it is usually correct; it misses many true positives.

\section{Common pitfalls (and how to avoid them)}
\begin{itemize}[leftmargin=*]
  \item \textbf{Using accuracy only:} for rare positives, accuracy can hide a useless model.
  \item \textbf{Forgetting the threshold:} changing $\tau$ changes precision/recall; do not treat it as fixed.
  \item \textbf{Ignoring the cost of errors:} choose metrics based on real harm (FN vs FP).
  \item \textbf{Comparing models on different splits:} use consistent splits or CV for fair comparison.
  \item \textbf{Not checking PR curves in imbalance:} PR-AUC can be more informative than ROC-AUC.
\end{itemize}

\section{Exit questions (with answers)}

\subsection*{Q1. Define TP, FP, TN, FN with a medical test.}
\textbf{Answer:} TP = sick predicted sick; FP = healthy predicted sick; FN = sick predicted healthy; TN = healthy predicted healthy.

\subsection*{Q2. Why is accuracy misleading on imbalanced datasets?}
\textbf{Answer:} a model can achieve high accuracy by always predicting the majority class while failing on the minority class.

\subsection*{Q3. When is recall more important than precision?}
\textbf{Answer:} when missing a positive is costly, e.g., disease screening or fraud detection.

\subsection*{Q4. Role of threshold?}
\textbf{Answer:} it converts model scores to labels and controls precision/recall trade-off.

\section{Summary}
Classification requires careful evaluation beyond accuracy. Confusion matrix metrics and threshold decisions are central, especially under imbalance.
\notesources{\AIMLUnitIVSources}

\end{document}
